% Five books, compared. Numbers from build/compare/numbers.tex (\lab{c-<slug>-...})
\input{../build/compare/numbers.tex}
\newcommand\cdir{../build/compare/}
\newcommand\bk[1]{\emph{\lab{c-#1-title}}}
\newcommand\cz[2]{\ifnum\lab{c-#1-sig-#2}=1 \textbf{\lab{c-#1-z-#2}}\else\lab{c-#1-z-#2}\fi}
% a row of the five books: #1 file prefix, #2 width each
\newcommand\fiverow[2]{%
  \noindent\begin{tabular}{@{}*{5}{c}@{}}
  \includegraphics[width=#2]{\cdir #1-the-trial.pdf} &
  \includegraphics[width=#2]{\cdir #1-bartleby.pdf} &
  \includegraphics[width=#2]{\cdir #1-yellow-wallpaper.pdf} &
  \includegraphics[width=#2]{\cdir #1-heart-of-darkness.pdf} &
  \includegraphics[width=#2]{\cdir #1-pride-and-prejudice.pdf}\\
  \fivenames
  \end{tabular}}
\newcommand\fivenames{{\figfont\scriptsize\itshape The Trial} & {\figfont\scriptsize\itshape Bartleby} &
  {\figfont\scriptsize\itshape Yellow Wallpaper} & {\figfont\scriptsize\itshape Heart of Darkness} &
  {\figfont\scriptsize\itshape Pride and Prejudice}}
\newcommand\stacked[2]{% #1 file prefix, #2 slug  — a full-width strip with its title
  \noindent{\figfont\scriptsize\itshape \lab{c-#2-title}}\\[0.5mm]
  \includegraphics[width=\textwidth]{\cdir #1-#2.pdf}\par\medskip}

\chapter{Five books}
\label{ch:compare}

A single drawing says little until there is another to set beside it. This
chapter walks four more books and lays them next to \emph{The Trial}. Each
was chosen to answer it in some way:

\begin{description}
\item[\bk{bartleby}] (\lab{c-bartleby-year}), Herman Melville's story of a
law-copyist who ``would prefer not to'', has often been read as Kafka before
Kafka: an office, a quiet refusal, a man undone without anyone quite meaning it.
\item[\bk{yellow-wallpaper}] (\lab{c-yellow-wallpaper-year}), Charlotte Perkins
Gilman's account of a woman confined to a room for her own good, written as
her private journal.
\item[\bk{heart-of-darkness}] (\lab{c-heart-of-darkness-year}), Joseph Conrad's
journey up a river, told aloud by Marlow on the deck of a boat at dusk.
\item[\bk{pride-and-prejudice}] (\lab{c-pride-and-prejudice-year}), Jane
Austen's comedy of manners: the contrast, a novel of drawing rooms, letters and
marriages, where nobody is put on trial except by gossip.
\end{description}

\noindent All four are in English and in the public domain. All five are
scored, walked and shuffled in exactly the same way.

\clearpage
\section{Five walks}

\begin{center}
\fiverow{shape}{22mm}

\bigskip
\fiverow{scale}{22mm}
\end{center}

\noindent Above, each book fills its own frame, so only the shape compares;
below, all five are drawn at one scale, so the size of the walk shows. Length
is the first and plainest difference. \bk{yellow-wallpaper} is a pocket of
\lab{c-yellow-wallpaper-sentences} sentences; \bk{pride-and-prejudice} runs to
\lab{c-pride-and-prejudice-sentences}.

\begin{table}[h]
\centering\figfont\footnotesize
\setlength\tabcolsep{4pt}
\begin{tabular}{@{}lrrrrr@{}}
\toprule
 & \emph{Trial} & \emph{Bartleby} & \emph{Wallpaper} & \emph{Darkness} & \emph{Pride}\\
\midrule
Sentences & \lab{c-the-trial-sentences} & \lab{c-bartleby-sentences} & \lab{c-yellow-wallpaper-sentences} & \lab{c-heart-of-darkness-sentences} & \lab{c-pride-and-prejudice-sentences}\\
Words per sentence & \lab{c-the-trial-stride} & \lab{c-bartleby-stride} & \lab{c-yellow-wallpaper-stride} & \lab{c-heart-of-darkness-stride} & \lab{c-pride-and-prejudice-stride}\\
Warm / cold / plain (\%) & \lab{c-the-trial-left}/\lab{c-the-trial-right}/\lab{c-the-trial-straight} & \lab{c-bartleby-left}/\lab{c-bartleby-right}/\lab{c-bartleby-straight} & \lab{c-yellow-wallpaper-left}/\lab{c-yellow-wallpaper-right}/\lab{c-yellow-wallpaper-straight} & \lab{c-heart-of-darkness-left}/\lab{c-heart-of-darkness-right}/\lab{c-heart-of-darkness-straight} & \lab{c-pride-and-prejudice-left}/\lab{c-pride-and-prejudice-right}/\lab{c-pride-and-prejudice-straight}\\
Lean (left minus right turns) & \lab{c-the-trial-net} & \lab{c-bartleby-net} & \lab{c-yellow-wallpaper-net} & \lab{c-heart-of-darkness-net} & \lab{c-pride-and-prejudice-net}\\
\quad without \emph{like} & \lab{c-the-trial-net-r} & \lab{c-bartleby-net-r} & \lab{c-yellow-wallpaper-net-r} & \lab{c-heart-of-darkness-net-r} & \lab{c-pride-and-prejudice-net-r}\\
Moods linger ($z$, same turn twice) & \cz{the-trial}{recoil} & \cz{bartleby}{recoil} & \cz{yellow-wallpaper}{recoil} & \cz{heart-of-darkness}{recoil} & \cz{pride-and-prejudice}{recoil}\\
Predictable turns ($z$, $H_2$) & \cz{the-trial}{entropy_h2} & \cz{bartleby}{entropy_h2} & \cz{yellow-wallpaper}{entropy_h2} & \cz{heart-of-darkness}{entropy_h2} & \cz{pride-and-prejudice}{entropy_h2}\\
Plain runs ($z$) & \cz{the-trial}{corridor_mean} & \cz{bartleby}{corridor_mean} & \cz{yellow-wallpaper}{corridor_mean} & \cz{heart-of-darkness}{corridor_mean} & \cz{pride-and-prejudice}{corridor_mean}\\
\bottomrule
\end{tabular}
\end{table}

\noindent \bk{pride-and-prejudice} leans warm by a wide margin, and stays warm
without \emph{like}. \bk{heart-of-darkness} and \bk{bartleby} lean cold, and
colder still without it. \bk{yellow-wallpaper} leans a little warm. Of the five,
\bk{the-trial} is the only one whose lean changes direction with that one word.
Its quietness is so even that a comparison can tip it.

The moods of \bk{pride-and-prejudice} linger even more strongly than those of
\bk{the-trial}: its turns are the most predictable of all five books. Austen's
chapters settle into a tone, a conversation of compliments or a scene of
mortification, and stay there. In the three shorter books, the same tendency is
present but not strong enough to tell from chance. They are short, and short
walks give chance more room, so this is an absence of evidence, not evidence
of absence.

\clearpage
\section{The vocabulary records}

\begin{center}
\vocabdisc{\cdir vocab-the-trial.png}{36mm}\hfill
\vocabdisc{\cdir vocab-bartleby.png}{36mm}\hfill
\vocabdisc{\cdir vocab-yellow-wallpaper.png}{36mm}\\[-1mm]
\makebox[36mm]{\figfont\scriptsize\itshape The Trial}\hfill\makebox[36mm]{\figfont\scriptsize\itshape Bartleby}\hfill\makebox[36mm]{\figfont\scriptsize\itshape The Yellow Wallpaper}\\[2mm]
\vocabdisc{\cdir vocab-heart-of-darkness.png}{36mm}\hspace{6mm}
\vocabdisc{\cdir vocab-pride-and-prejudice.png}{36mm}\\[-1mm]
\makebox[36mm]{\figfont\scriptsize\itshape Heart of Darkness}\hspace{6mm}\makebox[36mm]{\figfont\scriptsize\itshape Pride and Prejudice}
\end{center}

\noindent Each disc is the whole dictionary, \lab{vocab-words} words laid out
like the groove of a record: the commonest word at the centre, then outward,
ring after touching ring, to the rarest at the rim. Every word a book uses is a
black dot. All five records have a black core, since every book needs
\emph{the}, \emph{and} and \emph{was}, and they differ in how far the dots
scatter towards the rim.

Long books use more words simply because they are long, so the fair comparison
is a stretch of equal length. In any 6,000 running words, \bk{heart-of-darkness}
uses \lab{c-heart-of-darkness-types-6k} different dictionary words and
\bk{bartleby} \lab{c-bartleby-types-6k}. \bk{pride-and-prejudice} uses
\lab{c-pride-and-prejudice-types-6k}, \bk{yellow-wallpaper}
\lab{c-yellow-wallpaper-types-6k}, and \bk{the-trial} only
\lab{c-the-trial-types-6k}, the plainest of the five. Kafka (or at least this
English Kafka) does his strangest work with the most ordinary words. Half of
\bk{yellow-wallpaper}'s vocabulary lies within the commonest
\lab{c-yellow-wallpaper-median}\,\% of the dictionary, the tightest core of
all, which suits a journal written by someone forbidden to work. The rarest words
each book reaches for are small portraits: \lab{c-heart-of-darkness-rarest}
for Conrad; \lab{c-pride-and-prejudice-rarest} for Austen;
\lab{c-the-trial-rarest} for Kafka's translator.

\clearpage
\section{Tree rings}

\begin{center}
\fiverow{rings}{22mm}
\end{center}

\noindent Cut through a book like a tree trunk. The first page is the centre,
each chapter a ring whose area is its length, and each sentence a small radial
grain, long where the sentence uses long words and short where its words are
plain. (Books without chapters are cut into ten equal rings.) \bk{pride-and-prejudice}
shows the fine, even growth of \lab{c-pride-and-prejudice-rings} short chapters,
\bk{heart-of-darkness} three broad seasons. In \bk{the-trial} the arrest
forms a broad core, but the widest band of growth is chapter~7, the lawyer,
the manufacturer and the painter, as if the case grew fastest in the year of
consultations. The grain is lightest in \bk{yellow-wallpaper}, whose words
average \lab{c-yellow-wallpaper-wordlen} letters, against
\lab{c-pride-and-prejudice-wordlen} in \bk{pride-and-prejudice}.

\clearpage
\section{The loom}

\begin{center}
\fiverow{loom}{22mm}
\end{center}

\noindent Here the cast is woven. Each column is a warp thread for one of the
names a book uses most, people and places, in the order they first appear.
The book runs down the page, and a black stitch marks every stretch in which a
name occurs.

In \bk{the-trial} the first thread is a solid black column: K. is named in
every stretch of the book. Everyone else comes and goes: the warders
\lab{c-the-trial-pair}, whose stitches run together, then Leni, Titorelli and
Block, each in their own episode. \bk{pride-and-prejudice} is the most woven
cloth: Elizabeth is present in \lab{c-pride-and-prejudice-lead-share}\,\% of the
book, and Jane is stitched alongside her more than any other pair. In
\bk{heart-of-darkness} the most constant name is Kurtz, and he is present in only
\lab{c-heart-of-darkness-lead-share}\,\% of the book. The novel is organised
around a man who is mostly absent, talked about before he is met. In
\bk{yellow-wallpaper} there is almost no cloth at all, two names
(\lab{c-yellow-wallpaper-cast}) in a narrow strip. The narrator is alone with
the paper.

\noindent{\figfont\scriptsize Columns, left to right. \emph{Trial}: \lab{c-the-trial-cast}.
\emph{Bartleby}: \lab{c-bartleby-cast}. \emph{Wallpaper}: \lab{c-yellow-wallpaper-cast}.
\emph{Darkness}: \lab{c-heart-of-darkness-cast}. \emph{Pride}: \lab{c-pride-and-prejudice-cast}.}

\clearpage
\section{Echoes}

\stacked{echo}{the-trial}
\stacked{echo}{bartleby}
\stacked{echo}{yellow-wallpaper}
\stacked{echo}{heart-of-darkness}
\stacked{echo}{pride-and-prejudice}

\noindent Every word that recurs, other than the commonest words and names, is
drawn as a single arch over the length of the book, spanning its longest
absence: from the place the book last used it to the place where it comes back.
A low arch is a word the book never forgets for long; a high arch is a word
that disappears and returns from far away. The longest absences are telling. In
\bk{yellow-wallpaper} they belong to \lab{c-yellow-wallpaper-echo-longest}:
the paper itself and the light that moves on it. In \bk{heart-of-darkness},
\emph{horror} is among the words that vanish longest before they return, as it
does at the end. \bk{pride-and-prejudice} lets its words return soonest. Of all
the returns of its words, only \lab{c-pride-and-prejudice-echo-far}\,\% come
after more than a quarter of the book: its world renews its vocabulary as it
goes, from ball to letter to visit.

\clearpage
\section{Tides of speech}

\stacked{tide}{the-trial}
\stacked{tide}{bartleby}
\stacked{tide}{yellow-wallpaper}
\stacked{tide}{heart-of-darkness}
\stacked{tide}{pride-and-prejudice}

\noindent The black tide is the share of words spoken inside quotation marks,
read through each book from left to right and folded into three rows. Two
books are almost all one thing. \bk{heart-of-darkness} is a flood
(\lab{c-heart-of-darkness-speech}\,\% of its words are inside quotation marks)
because almost the whole novel is Marlow's voice, spoken aloud on the deck of
the Nellie. The few dry stretches are the unnamed listener who frames him.
\bk{yellow-wallpaper} is nearly dry (\lab{c-yellow-wallpaper-speech}\,\%): a
journal, kept in secret, where other people's voices are only remembered.
Between them, \bk{pride-and-prejudice} is half conversation
(\lab{c-pride-and-prejudice-speech}\,\%), and \bk{the-trial}
(\lab{c-the-trial-speech}\,\%) alternates long talking scenes with stretches
of narrative silence: the longest runs \lab{c-the-trial-silence} words without a
spoken word.

\clearpage
\section{Five portraits}

Put together, the instruments give each book a face.

\bk{the-trial} walks evenly, in quiet moods that linger, with the plainest
vocabulary of the five, around one man who is never absent. Its tangles sit in
the rooms where he looks for answers.

\bk{bartleby} is short and cool, and leans cold. Its cast is the office:
Turkey and Nippers, stitched together, and Bartleby himself. Its vocabulary is
rich for its length; a lawyer narrates it.

\bk{yellow-wallpaper} is the smallest walk, with the plainest words, almost no
other voices and almost no cast, and its longest echoes are the yellow paper
and the moonlight.

\bk{heart-of-darkness} is one long spoken voice, leaning cold, with the
richest vocabulary of the five, organised around a man who is mostly not there.

\bk{pride-and-prejudice} is the largest walk and the warmest, a densely woven
cloth of names around Elizabeth, half spoken, with moods that settle and stay.
